%% GENERATED by tools/gen_latex.py -- edit content/agents04.py instead.
%% Build:  latexmk -pdf agents04.tex     (or: pdflatex agents04.tex, twice)
\documentclass[aspectratio=169,11pt,t]{beamer}
\usepackage{itbpro}

\coursetitle{Ingatan: yang disimpan, yang dicari lagi, yang dibuang}
\coursesubtitle{Model tidak mengingat apa pun. Semua yang tampak seperti ingatan adalah keputusan rekayasa --- dan tiap keputusannya punya harga yang bisa dihitung.}
\coursekicker{Chapter 4}
\coursesource{Grootendorst \& Alammar, An Illustrated Guide to AI Agents --- bab 4}
\courseduration{3 jam (2 sesi)}
\coursepresenter{Hendri Karisma}{Instructor}
\coursebrand{AI for Professional \textperiodcentered\ ITB \textperiodcentered\ Ch 4}

\title{Ingatan: yang disimpan, yang dicari lagi, yang dibuang}
\author{Hendri Karisma}
\date{}

\begin{document}
\covertitle

\framekicker{Learning outcomes}
\begin{frame}[shrink=25]{Session Objectives}
\leadpar{By the end of this session, participants can:}
\begin{isteps}
\item \textbf{Menyebutkan tiga lapis ingatan} dan menentukan apa yang pantas disimpan di masing-masing.
\item \textbf{Membandingkan memangkas dan meringkas}, dan menyebutkan kerugian yang tidak terlihat dari yang kedua.
\item \textbf{Menjelaskan pengambilan sebagai ALAT}, bukan sebagai langkah praproses --- dan menyebutkan akibatnya pada jejak dan kutipan.
\item \textbf{Menyebutkan empat operasi rekayasa konteks} dan menerapkannya pada satu gelung yang terlalu panjang.
\item \textbf{Menyebutkan dua kelas serangan} yang muncul begitu agen menyimpan ingatan, dan penanganannya di batas alat.
\item \textbf{Memutuskan apa yang TIDAK boleh disimpan}, dengan alasan yang bisa dipertahankan di depan pemeriksa.
\end{isteps}

\vskip 6pt
\vskip 4pt
\reslink{Site}{Course home}{https://hendrikarisma.my.id/rs/ai-products-course/index.html}
\reslink{Book}{Maarten Grootendorst \& Jay Alammar, An Illustrated Guide to AI Agents (O'Reilly Media, Early Release)}{https://www.oreilly.com/library/view/an-illustrated-guide/9798341662681/}
\vskip 2pt
\end{frame}

\coursesection{01}{Tiga lapis, dan tidak satu pun gratis}{Yang tampak seperti ingatan selalu keputusan seseorang.}

\framekicker{Dasar}
\begin{frame}[shrink=25]{Yang disebut ingatan sebenarnya daftar pesan Anda}
{\fontsize{9.4}{12.6}\selectfont\color{fg2}Bab 2 sudah menyebutkannya dan bab ini dibangun di atasnya: \textbf{model tidak menyimpan apa pun antar panggilan.} Yang membuatnya tampak mengingat adalah kode Anda yang mengirim ulang riwayatnya.\par}\vskip 4pt
{\fontsize{9.4}{12.6}\selectfont\color{fg2}Itu berarti seluruh perilaku \textbackslash\{\}u201cmengingat\textbackslash\{\}u201d pada sistem Anda adalah \textbf{rekayasa} --- apa yang dimasukkan, apa yang dibuang, dan apa yang dicari kembali. Tidak ada yang otomatis, dan tidak ada yang gratis.\par}\vskip 4pt
\begin{iband}[signal]
Kabar baiknya: keadaan yang Anda pegang adalah keadaan yang bisa Anda \textbf{periksa, potong, uji, dan pertanggungjawabkan}. Ingatan yang ada di dalam bobot model tidak punya satu pun dari sifat itu.
\end{iband}
\end{frame}

\framekicker{Dasar}
\begin{frame}[shrink=25]{Tiga lapis, dengan umur dan harga yang berbeda}
\begin{center}
\includegraphics[width=0.92\linewidth,height=0.60\textheight,keepaspectratio]{figs/agents04-tiers.tex.pdf}
\end{center}
\vskip 3pt{\color{ink3}\fontsize{7.4}{9.4}\selectfont Semakin ke kanan, semakin murah disimpan dan semakin mahal diambil.\par}
{\fontsize{9.4}{12.6}\selectfont\color{fg2}\textbf{Ingatan kerja} hidup satu giliran: hasil alat mentah, yang biasanya besar dan hampir seluruhnya tidak perlu dibawa ke giliran berikutnya. \textbf{Jangka pendek} adalah percakapan ini. \textbf{Jangka panjang} ada di luar konteks dan hanya masuk ketika dicari.\par}\vskip 4pt
\end{frame}

\framekicker{Dasar}
\begin{frame}[shrink=25]{Pertanyaan yang memutuskan lapis mana}
\begin{itable}{>{\raggedright\arraybackslash}p{0.3572\textwidth}>{\raggedright\arraybackslash}p{0.2444\textwidth}>{\raggedright\arraybackslash}p{0.3384\textwidth}}
\thead{Pertanyaan} & \thead{Kalau ya} & \thead{Contoh} \\ \midrule
Dipakai lagi di giliran berikutnya? & Jangka pendek & Nomor rekening yang sedang dibahas \\
Dipakai lagi minggu depan? & Jangka panjang & Preferensi pengguna, keputusan yang pernah diambil \\
Hanya untuk menghasilkan satu angka? & Ingatan kerja & 1 843 baris transaksi $\rightarrow$ satu rasio \\
Bisa dihitung ulang kapan saja? & \textbf{Jangan disimpan} & Apa pun yang alat bisa berikan lagi \\
\end{itable}
\begin{iband}[signal]
Baris terakhir yang paling sering dilanggar. Menyimpan sesuatu yang bisa diambil ulang berarti menukar \hilite{biaya token yang pasti dengan risiko data basi}
\end{iband}
\end{frame}

\coursesection{02}{Jangka pendek: memangkas atau meringkas}{Dua cara, dua kerugian --- dan hanya satu yang kelihatan.}

\framekicker{Jangka pendek}
\begin{frame}[shrink=25]{Memangkas: kehilangan yang jujur}
{\fontsize{9.4}{12.6}\selectfont\color{fg2}Cara paling sederhana menjaga konteks tetap muat: simpan N giliran terakhir, buang sisanya. Murah, pasti, dan mudah diuji.\par}\vskip 4pt
{\fontsize{9.4}{12.6}\selectfont\color{fg2}Kerugiannya juga pasti, dan itu kekuatannya: Anda tahu persis apa yang hilang. Masalahnya muncul ketika yang hilang adalah \textbf{batasan yang disebut di awal} --- \textbackslash\{\}u201cjangan hubungi nasabah lewat surel\textbackslash\{\}u201d yang diucapkan pada giliran kedua tidak ada lagi pada giliran kedua puluh.\par}\vskip 4pt
\begin{iband}[signal]
Bentuk kegagalannya khas dan mudah dikenali: agen bekerja benar selama sepuluh giliran, lalu mulai \textbf{melanggar aturan yang sudah tidak ada di konteksnya}. Itu bukan model yang lupa; itu jendela yang menggeser.
\end{iband}
\end{frame}

\framekicker{Jangka pendek}
\begin{frame}[shrink=25]{Meringkas: kehilangan yang tidak terlihat}
\begin{center}
\includegraphics[width=0.97\linewidth,height=0.56\textheight,keepaspectratio]{figs/agents04-decay.tex.pdf}
\end{center}
\vskip 3pt{\color{ink3}\fontsize{7.4}{9.4}\selectfont Peringkasan yang diulang, dihitung pemangkatannya. Langkahi: asli, dua kali pertama, lalu enam kali.\par}
{\fontsize{9.4}{12.6}\selectfont\color{fg2}Meringkas terasa lebih baik daripada memangkas --- tidak ada yang dibuang, semuanya \textbackslash\{\}u201cmasih ada\textbackslash\{\}u201d. Tetapi peringkasan berikutnya bekerja atas \textbf{ringkasan}, bukan atas aslinya, sehingga apa pun yang hilang tiap kali dipangkatkan.\par}\vskip 4pt
\end{frame}

\framekicker{Jangka pendek}
\begin{frame}[shrink=25]{Yang paling dulu hilang dari sebuah ringkasan}
\begin{minipage}[t]{0.494\textwidth}
\begin{icardbad}
\cardh{Angka dan pengenal}
\textbackslash\{\}u201cRasio 0,82\textbackslash\{\}u201d jadi \textbackslash\{\}u201crasionya rendah\textbackslash\{\}u201d. Setelah itu tidak ada cara memulihkannya, dan tidak ada yang tahu itu terjadi.
\end{icardbad}
\end{minipage}%
\hfill
\begin{minipage}[t]{0.494\textwidth}
\begin{icardbad}
\cardh{Batasan yang bersyarat}
\textbackslash\{\}u201cBoleh, KECUALI kalau nasabahnya di bawah 21\textbackslash\{\}u201d cenderung jadi \textbackslash\{\}u201cboleh\textbackslash\{\}u201d.
\end{icardbad}
\end{minipage}%


\vskip 4pt

\begin{minipage}[t]{0.494\textwidth}
\begin{icard}
\cardh{Yang belum diputuskan}
Ringkasan cenderung menutup pertanyaan yang masih terbuka, sebab teks yang rapi terbaca lebih baik.
\end{icard}
\end{minipage}%
\hfill
\begin{minipage}[t]{0.494\textwidth}
\begin{icard}
\cardh{Urutan waktu}
Mana yang lebih dulu sering hilang, dan itu penting justru saat menelusuri kesalahan.
\end{icard}
\end{minipage}%
{\fontsize{9.4}{12.6}\selectfont\color{fg2}Perhatikan polanya: yang hilang duluan justru yang \textbf{paling menentukan keputusan}. Prosa yang menjelaskan bertahan; angka dan pengecualian tidak.\par}\vskip 4pt
\end{frame}

\framekicker{Jangka pendek}
\begin{frame}[shrink=25]{Bentuk campuran yang biasanya menang}
\begin{center}

\includegraphics[width=0.92\linewidth,height=0.60\textheight,keepaspectratio]{figs/agents04-hybrid.tex.pdf}
\end{center}
\vskip 3pt{\color{ink3}\fontsize{7.4}{9.4}\selectfont Yang tetap, yang diringkas, dan yang utuh --- masing-masing punya alasan.\par}
\begin{isteps}
\item \{'h': 'Perintah sistem tidak pernah diringkas', 'p': 'Ia bagian tetap, dan ia yang membuat singgahan prompt bekerja.'\}
\item \{'h': 'Giliran pertama disimpan utuh', 'p': 'Di situ tujuan dan batasan disebut. Inilah yang paling sering hilang dan paling mahal hilangnya.'\}
\item \{'h': 'Bagian tengah diringkas, satu kali', 'p': 'Ringkas dari \textbf{aslinya} kalau masih ada, bukan dari ringkasan sebelumnya. Itu yang mematahkan pemangkatan di gambar tadi.'\}
\item \{'h': 'Beberapa giliran terakhir utuh', 'p': 'Di situ konteks langsung berada, dan di situ ketepatan paling terasa.'\}
\end{isteps}
\end{frame}

\framekicker{Jangka pendek}
\begin{frame}[shrink=25]{Simpan aslinya, ringkas dari sana}
{\fontsize{9.4}{12.6}\selectfont\color{fg2}Satu perubahan kecil yang menghapus sebagian besar masalah di gambar tadi: \textbf{simpan percakapan aslinya di luar konteks}, dan buat tiap ringkasan dari yang asli, bukan dari ringkasan sebelumnya.\par}\vskip 4pt
{\fontsize{9.4}{12.6}\selectfont\color{fg2}Penyimpanannya murah --- teks di basis data, bukan token di konteks. Yang mahal adalah konteks, dan itu tidak bertambah.\par}\vskip 4pt
\begin{iband}[signal]
Dengan begitu kerugiannya berhenti berlipat: ia jadi kerugian \textbf{satu kali} dari sumber lengkap, bukan kerugian berulang dari salinan yang makin tipis.
\end{iband}
\end{frame}

\coursesection{03}{Jangka panjang: mengambil, bukan mengingat}{Dan kenapa pengambilan itu alat, bukan langkah praproses.}

\framekicker{Jangka panjang}
\begin{frame}[shrink=25]{Bentuk dasarnya, dan satu kalimat yang menentukan segalanya}
\begin{center}
\includegraphics[width=0.92\linewidth,height=0.60\textheight,keepaspectratio]{figs/agents04-rag.tex.pdf}
\end{center}
\vskip 3pt{\color{ink3}\fontsize{7.4}{9.4}\selectfont Pertanyaan jadi vektor, dicari yang mirip, potongan teratas masuk konteks.\par}
{\fontsize{9.4}{12.6}\selectfont\color{fg2}Mekanismenya sudah dibahas di \textbf{Bab 14} kelas ini --- penyematan, kemiripan, indeks. Yang penting di sini satu kalimat: \textbf{model hanya bisa menjawab dari potongan yang terambil.} Apa pun yang tidak terambil, untuk keperluan jawaban itu, tidak ada.\par}\vskip 4pt
\begin{iband}[signal]
Karena itu mutu sistem semacam ini hampir seluruhnya ditentukan oleh \textbf{mutu pengambilannya}, bukan oleh kepintaran modelnya --- dan di situlah waktu perbaikan sebaiknya dihabiskan.
\end{iband}
\end{frame}

\framekicker{Jangka panjang}
\begin{frame}[shrink=25]{Pengambilan sebagai alat mengubah tiga hal}
\begin{center}
\includegraphics[width=0.92\linewidth,height=0.60\textheight,keepaspectratio]{figs/agents04-agentic.tex.pdf}
\end{center}
\vskip 3pt{\color{ink3}\fontsize{7.4}{9.4}\selectfont Agen memutuskan apa yang dicari, membaca hasilnya, dan boleh mencari lagi.\par}
\begin{minipage}[t]{0.485\textwidth}
{\fontsize{9.4}{12.6}\selectfont\color{fg2}\textbf{Praproses}: satu pencarian, dengan kata-kata pengguna apa adanya, sebelum model melihat apa pun.\par}\vskip 4pt
{\fontsize{9.4}{12.6}\selectfont\color{fg2}Kalau kata pengguna buruk, hasilnya buruk, dan tidak ada kesempatan kedua.\par}\vskip 4pt

\end{minipage}%
\hfill
\begin{minipage}[t]{0.485\textwidth}
{\fontsize{9.4}{12.6}\selectfont\color{fg2}\textbf{Alat}: model menyusun kata pencariannya sendiri, membaca hasilnya, dan boleh mencari lagi dengan istilah berbeda.\par}\vskip 4pt
{\fontsize{9.4}{12.6}\selectfont\color{fg2}Lebih mahal --- tiap pencarian satu giliran --- dan jauh lebih tahan terhadap pertanyaan yang dirumuskan buruk.\par}\vskip 4pt

\end{minipage}%
\begin{iband}[signal]
Dan yang paling penting untuk sistem yang diawasi: sebagai alat, tiap pengambilan \textbf{tercatat di jejak} --- apa yang dicari, apa yang kembali, dan potongan mana yang dipakai.
\end{iband}
\end{frame}

\framekicker{Jangka panjang}
\begin{frame}[shrink=25]{Kutipan bukan hiasan; ia satu-satunya cara memeriksa}
{\fontsize{9.4}{12.6}\selectfont\color{fg2}Jawaban yang menyebut sumbernya bisa diperiksa dalam hitungan detik. Jawaban yang tidak, hanya bisa dipercaya atau tidak.\par}\vskip 4pt
\begin{isteps}
\item \{'h': 'Kutip potongan, bukan dokumen', 'p': '\textbackslash\{\}\textbackslash\{\}u201cKebijakan kredit\textbackslash\{\}\textbackslash\{\}u201d tidak bisa diperiksa; \textbackslash\{\}\textbackslash\{\}u201cklausul CP-04, paragraf 2\textbackslash\{\}\textbackslash\{\}u201d bisa.'\}
\item \{'h': 'Periksa kutipannya dengan kode', 'p': 'Kalau pengenal yang dikutip tidak ada di hasil pengambilan, itu karangan --- dan kode Anda bisa mengetahuinya tanpa model.'\}
\item \{'h': 'Simpan potongannya di jejak', 'p': 'Agar setahun kemudian pertanyaan \textbackslash\{\}\textbackslash\{\}u201catas dasar apa\textbackslash\{\}\textbackslash\{\}u201d punya jawaban, bukan rekonstruksi.'\}
\end{isteps}
\begin{iband}[signal]
Demo kredit UMKM melakukan persis ini: tiap rekomendasi membawa klausul yang dipakai, dicetak penuh, dan uji akan gagal kalau klausul yang dikutip tidak ada.
\end{iband}
\end{frame}

\framekicker{Jangka panjang}
\begin{frame}[shrink=25]{Apa yang pantas jadi ingatan jangka panjang}
\begin{minipage}[t]{0.3253\textwidth}
\begin{icardgood}
\cardh{Fakta yang stabil}
Preferensi, keputusan yang pernah diambil beserta alasannya, hasil yang mahal dihitung ulang.
\end{icardgood}
\end{minipage}%
\hfill
\begin{minipage}[t]{0.3253\textwidth}
\begin{icard}
\cardh{Perlu kedaluwarsa}
Apa pun yang bisa berubah di sistem lain. Simpan pengenalnya, ambil nilainya saat dipakai.
\end{icard}
\end{minipage}%
\hfill
\begin{minipage}[t]{0.3253\textwidth}
\begin{icardbad}
\cardh{Jangan sama sekali}
Data pribadi yang tidak dibutuhkan agen, kredensial, dan apa pun yang retensinya diatur --- kecuali retensi itu ikut diterapkan.
\end{icardbad}
\end{minipage}%
{\fontsize{9.4}{12.6}\selectfont\color{fg2}Kartu ketiga bukan nasihat kehati-hatian umum. Ingatan adalah \textbf{basis data baru} yang sering lahir tanpa pemilik, tanpa kebijakan retensi, dan tanpa seorang pun menyadari ia sudah berisi data pribadi.\par}\vskip 4pt
\end{frame}

\framekicker{Jangka panjang}
\begin{frame}[shrink=25]{Di mana teks dipotong menentukan apa yang bisa ditemukan}
\begin{center}
\includegraphics[width=0.97\linewidth,height=0.56\textheight,keepaspectratio]{figs/agents04-chunk.tex.pdf}
\end{center}
\vskip 3pt{\color{ink3}\fontsize{7.4}{9.4}\selectfont Satu kalimat, tiga cara memotong. Langkahi: ukuran tetap, tumpang tindih, lalu menurut struktur.\par}
{\fontsize{9.4}{12.6}\selectfont\color{fg2}Pemotongan biasanya digambar sebagai deretan kotak seukuran, yang memperlihatkan mekanismenya dan menyembunyikan satu-satunya hal yang penting: \textbf{kalimat yang maknanya melintasi batas potongan hilang dari kedua sisinya.}\par}\vskip 4pt
\end{frame}

\framekicker{Jangka panjang}
\begin{frame}[shrink=25]{Memilih cara memotong, tanpa mencoba semuanya}
\begin{itable}{>{\raggedright\arraybackslash}p{0.282\textwidth}>{\raggedright\arraybackslash}p{0.282\textwidth}>{\raggedright\arraybackslash}p{0.376\textwidth}}
\thead{Bentuk dokumen} & \thead{Cara yang cocok} & \thead{Alasannya} \\ \midrule
Kebijakan, peraturan, kontrak & Menurut struktur & Klausul memang unit maknanya, dan nomornya jadi kutipan \\
Manual, dokumentasi & Menurut judul bagian & Satu bagian menjawab satu pertanyaan \\
Percakapan, tiket & Per giliran atau per pesan & Batas alaminya sudah ada \\
Prosa panjang tanpa struktur & Tetap + tumpang tindih & Tidak ada batas alami; tumpang tindih membeli jaminan \\
Tabel dan lembar kerja & \textbf{Jangan dipotong sebagai teks} & Jadikan alat kueri, bukan potongan untuk dicari kemiripannya \\
\end{itable}
\begin{iband}[signal]
Baris terakhir yang paling sering dilanggar dan paling mahal: \hilite{tabel yang diubah jadi potongan teks kehilangan justru sifat yang membuatnya berguna}, yaitu bisa dihitung.
\end{iband}
\end{frame}

\framekicker{Jangka panjang}
\begin{frame}[shrink=25]{Kemiripan bukan relevansi}
{\fontsize{9.4}{12.6}\selectfont\color{fg2}Pencarian vektor menemukan yang \textbf{mirip secara makna}. Itu bukan hal yang sama dengan yang \textbf{menjawab pertanyaan}, dan bedanya muncul dengan cara yang khas.\par}\vskip 4pt
\begin{minipage}[t]{0.3253\textwidth}
\begin{icardbad}
\cardh{Parafrase pertanyaan}
Potongan yang mengulang pertanyaan dengan kata lain skornya tinggi, isinya tidak menjawab apa pun.
\end{icardbad}
\end{minipage}%
\hfill
\begin{minipage}[t]{0.3253\textwidth}
\begin{icardbad}
\cardh{Angka dan kode}
Kemiripan makna buruk untuk nomor klausul, kode produk, dan pengenal --- di situ pencarian kata kunci menang telak.
\end{icardbad}
\end{minipage}%
\hfill
\begin{minipage}[t]{0.3253\textwidth}
\begin{icardbad}
\cardh{Negasi}
``boleh'' dan ``tidak boleh'' sangat mirip secara vektor, dan berlawanan artinya.
\end{icardbad}
\end{minipage}%
\begin{iband}[signal]
Obat yang paling sering cukup: \textbf{gabungkan pencarian kata kunci dengan pencarian vektor}, lalu urutkan ulang hasil gabungannya. Ini rekayasa pencarian biasa, dan ia mengalahkan penggantian model.
\end{iband}
\end{frame}

\framekicker{Jangka panjang}
\begin{frame}[shrink=25]{Ambil banyak, saring ketat}
\begin{isteps}
\item \{'h': 'Ambil lebih banyak daripada yang akan dipakai', 'p': 'Pencarian awal murah dan tidak perlu tepat --- tugasnya memastikan jawabannya \textbf{ada di dalam} kumpulan yang diambil.'\}
\item \{'h': 'Urutkan ulang dengan model yang lebih teliti', 'p': 'Membaca pertanyaan dan potongan bersama-sama, bukan membandingkan dua vektor. Lebih mahal per potongan, dan hanya dijalankan pada yang sedikit.'\}
\item \{'h': 'Masukkan hanya yang teratas', 'p': 'Bab 2: konteks yang penuh menurunkan ketepatan. Yang tidak masuk bukan kerugian, ia penghematan ganda.'\}
\end{isteps}
\begin{iband}[signal]
Pola ini --- \textbf{jaring lebar lalu saringan halus} --- sama dengan deteksi objek di Bab 12: usulan banyak, lalu penyaringan. Bentuk yang sama muncul lagi karena alasan yang sama.
\end{iband}
\end{frame}

\coursesection{04}{Rekayasa konteks}{Empat operasi, dan konteks sebagai spesifikasi.}

\framekicker{Rekayasa konteks}
\begin{frame}[shrink=25]{Empat operasi, dan itu saja}
\begin{itable}{>{\raggedright\arraybackslash}p{0.188\textwidth}>{\raggedright\arraybackslash}p{0.376\textwidth}>{\raggedright\arraybackslash}p{0.376\textwidth}}
\thead{Operasi} & \thead{Pertanyaannya} & \thead{Kesalahan yang biasa} \\ \midrule
Melacak & Apa yang sudah terjadi dan di mana disimpan? & Tidak ada yang disimpan di luar konteks, jadi tidak ada yang bisa diambil ulang \\
Memilih & Apa yang masuk ke giliran ini? & Memasukkan semua yang muat, sebab \textbackslash\{\}u201cmasih ada ruang\textbackslash\{\}u201d \\
Memampatkan & Bagaimana memperkecil tanpa kehilangan yang penting? & Meringkas ringkasan --- lihat gambar tadi \\
Mengurutkan & Apa di awal, apa di akhir? & Menaruh aturan penting di tengah tumpukan hasil alat \\
\end{itable}
\begin{iband}[signal]
Keempatnya keputusan \textbf{kode}, bukan keputusan model. Kalau tidak ada yang menuliskannya, keempatnya tetap terjadi --- hanya saja secara kebetulan.
\end{iband}
\end{frame}

\framekicker{Rekayasa konteks}
\begin{frame}[shrink=25]{Memilih: yang muat bukan berarti pantas masuk}
{\fontsize{9.4}{12.6}\selectfont\color{fg2}Godaan terbesar dari jendela konteks yang besar adalah memasukkan semuanya. Bab 2 sudah menyebut sebabnya itu buruk: ketepatan menemukan informasi \textbf{menurun} pada konteks yang penuh, terutama di bagian tengah.\par}\vskip 4pt
{\fontsize{9.4}{12.6}\selectfont\color{fg2}Jadi memilih bukan penghematan biaya semata --- ia \textbf{menaikkan mutu jawaban}. Sepuluh potongan yang tepat mengalahkan seratus potongan yang mengandung sepuluh itu.\par}\vskip 4pt
\begin{iband}[signal]
Uji yang mudah dan jarang dilakukan: jalankan kumpulan uji Anda dengan potongan yang diambil dikurangi separuh. Kalau angkanya tidak turun, \hilite{separuh konteks Anda selama ini hanya biaya}
\end{iband}
\end{frame}

\framekicker{Rekayasa konteks}
\begin{frame}[shrink=25]{Mengurutkan: awal dan akhir adalah tempat istimewa}
\begin{minipage}[t]{0.485\textwidth}
{\fontsize{9.4}{12.6}\selectfont\color{fg2}\textbf{Taruh di awal}\par}\vskip 4pt
\begin{ibullets}
\item Peran dan batasan
\item Definisi yang dipakai sepanjang percakapan
\item Skema alat
\end{ibullets}

\end{minipage}%
\hfill
\begin{minipage}[t]{0.485\textwidth}
{\fontsize{9.4}{12.6}\selectfont\color{fg2}\textbf{Taruh di akhir}\par}\vskip 4pt
\begin{ibullets}
\item Permintaan giliran ini
\item Hasil alat yang baru saja kembali
\item Pengingat format keluaran
\end{ibullets}

\end{minipage}%
{\fontsize{9.4}{12.6}\selectfont\color{fg2}Dan konsekuensi yang sering terlewat: menaruh pengingat format di \textbf{akhir} bekerja jauh lebih baik daripada menaruhnya di perintah sistem, ketika percakapannya sudah panjang.\par}\vskip 4pt
\end{frame}

\framekicker{Rekayasa konteks}
\begin{frame}[shrink=25]{Konteks adalah spesifikasinya}
{\fontsize{9.4}{12.6}\selectfont\color{fg2}Pada sistem berbasis model, yang menentukan perilaku bukan kode yang Anda tulis, melainkan \textbf{konteks yang Anda susun}. Itu spesifikasi sistemnya, dan ia dirakit ulang tiap giliran.\par}\vskip 4pt
\begin{isteps}
\item \{'h': 'Kalau ia dirakit, ia bisa diuji', 'p': 'Rakit konteks untuk kasus uji, dan periksa isinya sebelum dikirim. Itu uji unit biasa, dan tidak butuh model.'\}
\item \{'h': 'Kalau ia spesifikasi, ia pantas ditinjau', 'p': 'Perubahan pada penyusun konteks pantas dibaca orang lain, seperti perubahan kode.'\}
\item \{'h': 'Kalau ia berubah tiap giliran, ia pantas dicatat', 'p': 'Jejak harus memuat konteks yang dipakai, bukan hanya jawabannya.'\}
\end{isteps}
\begin{iband}[signal]
Cara paling cepat memperbaiki agen yang mengecewakan hampir selalu \textbf{memperbaiki apa yang masuk ke konteksnya}, bukan mengganti modelnya.
\end{iband}
\end{frame}

\framekicker{Rekayasa konteks}
\begin{frame}[shrink=25]{Ingatan kerja: jangan biarkan hasil alat menetap}
{\fontsize{9.4}{12.6}\selectfont\color{fg2}Hasil alat sering besar dan hampir seluruhnya sekali pakai. Membiarkannya menetap di riwayat berarti membayarnya ulang di tiap giliran sesudahnya --- Bab 2 menghitung berapa lipat itu.\par}\vskip 4pt
\begin{minipage}[t]{0.485\textwidth}
{\fontsize{9.4}{12.6}\selectfont\color{fg2}\textbf{Pola yang mahal}\par}\vskip 4pt
{\fontsize{9.4}{12.6}\selectfont\color{fg2}Alat mengembalikan 1 843 baris transaksi; semuanya masuk riwayat; enam giliran berikutnya membayarnya lagi dan lagi.\par}\vskip 4pt

\end{minipage}%
\hfill
\begin{minipage}[t]{0.485\textwidth}
{\fontsize{9.4}{12.6}\selectfont\color{fg2}\textbf{Pola yang murah}\par}\vskip 4pt
{\fontsize{9.4}{12.6}\selectfont\color{fg2}Alat menghitung di sisinya sendiri dan mengembalikan rasio yang diminta. Baris mentahnya disimpan di luar konteks, dengan pengenal, kalau-kalau perlu.\par}\vskip 4pt

\end{minipage}%
\begin{iband}[signal]
Aturan yang mudah diingat: \textbf{alat mengembalikan jawaban, bukan bahan.} Kalau agen harus menghitung sendiri dari data mentah, itu tanda pekerjaan berada di tempat yang salah.
\end{iband}
\end{frame}

\framekicker{Rekayasa konteks}
\begin{frame}[shrink=25]{Siapa yang memutuskan sesuatu layak diingat}
{\fontsize{9.4}{12.6}\selectfont\color{fg2}Ada tiga cara, dan pilihannya menentukan seberapa besar ingatan Anda tumbuh serta seberapa bisa dipercaya isinya.\par}\vskip 4pt
\begin{itable}{>{\raggedright\arraybackslash}p{0.2068\textwidth}>{\raggedright\arraybackslash}p{0.376\textwidth}>{\raggedright\arraybackslash}p{0.3572\textwidth}}
\thead{Siapa} & \thead{Cara kerjanya} & \thead{Pertukarannya} \\ \midrule
Aturan kode & Simpan hal yang jenisnya sudah ditentukan & Bisa diramalkan dan diuji; melewatkan yang tak terduga \\
Model & Agen memanggil alat ``ingat ini'' & Menangkap yang tak terduga; tumbuh cepat dan bisa salah \\
Manusia & Pengguna menandai sesuatu untuk diingat & Paling tepat, paling jarang dipakai \\
\end{itable}
\begin{iband}[signal]
Yang paling sering berhasil: \textbf{kode untuk yang wajib, model untuk yang opsional, dan batas jumlah} --- supaya ``ingat ini'' tidak berubah jadi tempat pembuangan.
\end{iband}
\end{frame}

\framekicker{Rekayasa konteks}
\begin{frame}[shrink=25]{Melupakan adalah fitur, bukan kekurangan}
{\fontsize{9.4}{12.6}\selectfont\color{fg2}Ingatan yang hanya bertambah akan berubah jadi tumpukan yang isinya saling bertentangan: preferensi lama dan baru, keputusan yang sudah dicabut, fakta yang sudah berubah di sistem lain.\par}\vskip 4pt
\begin{isteps}
\item \{'h': 'Beri tanggal pada semuanya', 'p': 'Tanpa waktu, tidak ada cara memilih antara dua ingatan yang bertentangan.'\}
\item \{'h': 'Perbarui, jangan menumpuk', 'p': 'Ingatan baru tentang hal yang sama harus \textbf{mengganti} yang lama, bukan berdiri di sebelahnya.'\}
\item \{'h': 'Kedaluwarsakan yang bisa basi', 'p': 'Apa pun yang sumbernya sistem lain punya umur. Simpan pengenalnya, ambil nilainya.'\}
\item \{'h': 'Sediakan penghapusan yang sungguhan', 'p': 'Termasuk dari indeks pencarian. Ingatan yang terhapus dari basis data tapi masih ada di indeks belum terhapus.'\}
\end{isteps}
\end{frame}

\framekicker{Rekayasa konteks}
\begin{frame}[shrink=25]{Ingatan bersama antar agen menciptakan masalah baru}
{\fontsize{9.4}{12.6}\selectfont\color{fg2}Begitu ada lebih dari satu agen, muncul pertanyaan yang tidak ada pada satu agen: \textbf{konteks siapa yang dilihat siapa?} Bab 8 membahasnya penuh; di sini cukup bentuk masalahnya.\par}\vskip 4pt
\begin{minipage}[t]{0.3253\textwidth}
\begin{icard}
\cardh{Semuanya melihat semuanya}
Paling sederhana, dan konteksnya membengkak paling cepat. Tiap agen membayar percakapan agen lain.
\end{icard}
\end{minipage}%
\hfill
\begin{minipage}[t]{0.3253\textwidth}
\begin{icardaccent}
\cardh{Hanya lewat pesan}
Tiap agen punya konteks sendiri dan bertukar ringkasan. Lebih murah, dan kehilangan detail di tiap penyerahan.
\end{icardaccent}
\end{minipage}%
\hfill
\begin{minipage}[t]{0.3253\textwidth}
\begin{icard}
\cardh{Papan bersama}
Keadaan disimpan di luar, tiap agen membaca yang perlu. Paling bisa diskalakan, paling banyak kodenya.
\end{icard}
\end{minipage}%
\begin{iband}[signal]
Perhatikan bahwa ketiganya soal \textbf{ingatan}, bukan soal kecerdasan. Sebagian besar kesulitan sistem banyak-agen ternyata kesulitan konteks.
\end{iband}
\end{frame}

\framekicker{Ongkos}
\begin{frame}[shrink=25]{Ingatan tidak gratis, dan ongkosnya bukan hanya token}
\begin{itable}{>{\raggedright\arraybackslash}p{0.2256\textwidth}>{\raggedright\arraybackslash}p{0.3196\textwidth}>{\raggedright\arraybackslash}p{0.3948\textwidth}}
\thead{Ongkos} & \thead{Kapan terasa} & \thead{Yang sering terlewat} \\ \midrule
Penyematan & Saat menulis dan saat mencari & Menyematkan ulang seluruh korpus setiap ganti model penyemat \\
Penyimpanan & Tumbuh pelan & Indeks vektor jauh lebih besar daripada teksnya \\
Waktu pengambilan & Tiap giliran & Menambah satu perjalanan pulang-pergi ke tiap langkah \\
Token konteks & Tiap giliran sesudahnya & Potongan yang diambil dibayar ulang seperti riwayat \\
\textbf{Kewajiban hukum} & Saat diperiksa & Basis data baru berisi data orang, sering tanpa pemilik \\
\end{itable}
{\fontsize{9.4}{12.6}\selectfont\color{fg2}Baris terakhir bukan ongkos teknis, dan justru itu yang paling sering tidak dianggarkan siapa pun.\par}\vskip 4pt
\end{frame}

\framekicker{Ongkos}
\begin{frame}[shrink=25]{Kapan tidak perlu ingatan sama sekali}
\begin{minipage}[t]{0.494\textwidth}
\begin{icardgood}
\cardh{Tugas satu kali yang berdiri sendiri}
Menilai satu pengajuan, meringkas satu dokumen. Semua yang dibutuhkan bisa diambil dengan alat saat itu juga.
\end{icardgood}
\end{minipage}%
\hfill
\begin{minipage}[t]{0.494\textwidth}
\begin{icardgood}
\cardh{Yang bisa diambil ulang kapan saja}
Kalau alat bisa memberikannya lagi, menyimpannya hanya menukar biaya pasti dengan risiko basi.
\end{icardgood}
\end{minipage}%
{\fontsize{9.4}{12.6}\selectfont\color{fg2}Demo kredit UMKM tidak punya ingatan lintas percakapan sama sekali, dan itu \textbf{keputusan}, bukan kelalaian: tiap penilaian berdiri sendiri, semua bahannya diambil lewat alat, dan tidak ada satu pun data nasabah yang menetap di luar sistem intinya.\par}\vskip 4pt
\begin{iband}[signal]
Pertanyaan yang pantas ditanyakan sebelum memasang ingatan: \textbf{kegagalan mana yang sedang saya perbaiki?} Kalau tidak ada jawabannya, yang sedang dipasang adalah basis data baru tanpa alasan.
\end{iband}
\end{frame}

\coursesection{05}{Ingatan sebagai permukaan serangan}{Dua hal yang baru muncul begitu sistem mulai mengingat.}

\framekicker{Risiko}
\begin{frame}[shrink=25]{Perintah yang menumpang di dalam data yang dibaca}
\begin{center}
\includegraphics[width=0.92\linewidth,height=0.60\textheight,keepaspectratio]{figs/agents04-poison.tex.pdf}
\end{center}
\vskip 3pt{\color{ink3}\fontsize{7.4}{9.4}\selectfont Hasil alat masuk ke konteks yang sama dengan perintah --- dan terlihat sama.\par}
{\fontsize{9.4}{12.6}\selectfont\color{fg2}Agen membaca dokumen, tiket, atau halaman. Isinya bisa memuat kalimat yang ditujukan kepada \textbf{agennya}, bukan kepada pembacanya. Selama hasil alat masuk ke konteks yang sama dengan perintah, keduanya tidak terbedakan.\par}\vskip 4pt
\begin{iband}[signal]
Penanganannya \textbf{bukan} di prompt. Ia di batas alat: apa yang boleh dipanggil sesudah membaca sesuatu yang tidak tepercaya, dan apakah alat tulis termasuk di dalamnya.
\end{iband}
\end{frame}

\framekicker{Risiko}
\begin{frame}[shrink=25]{Dan ingatan membuatnya bertahan sampai besok}
{\fontsize{9.4}{12.6}\selectfont\color{fg2}Tanpa ingatan, serangan semacam itu berakhir ketika percakapannya berakhir. Dengan ingatan jangka panjang, kalimat yang disisipkan bisa \textbf{tersimpan} --- dan muncul lagi pada percakapan berikutnya, dengan pengguna yang berbeda.\par}\vskip 4pt
\begin{isteps}
\item \{'h': 'Jangan simpan teks mentah dari sumber tidak tepercaya', 'p': 'Simpan fakta yang sudah diekstrak dan divalidasi, bukan paragrafnya.'\}
\item \{'h': 'Beri asal-usul pada tiap ingatan', 'p': 'Dari percakapan mana, dari dokumen mana, ditulis kapan. Tanpa itu, tidak ada cara mencabutnya.'\}
\item \{'h': 'Sediakan cara menghapus', 'p': 'Ingatan yang tidak bisa dihapus adalah kewajiban hukum yang tidak bisa dipenuhi.'\}
\end{isteps}
\end{frame}

\framekicker{Risiko}
\begin{frame}[shrink=25]{Ingatan bersama antar pengguna adalah keputusan, bukan fitur}
{\fontsize{9.4}{12.6}\selectfont\color{fg2}Ingatan yang dibagi antar pengguna terdengar efisien: sistem belajar sekali, semua orang menikmatinya. Ia juga jalur kebocoran yang paling langsung.\par}\vskip 4pt
\begin{minipage}[t]{0.494\textwidth}
\begin{icardgood}
\cardh{Per pengguna, bawaan}
Ingatan diberi batas pemilik, dan pencarian selalu tersaring menurut pemiliknya --- di \textbf{kode}, bukan di prompt.
\end{icardgood}
\end{minipage}%
\hfill
\begin{minipage}[t]{0.494\textwidth}
\begin{icard}
\cardh{Bersama, kalau memang perlu}
Hanya untuk pengetahuan yang memang milik organisasi: kebijakan, prosedur, definisi. Bukan untuk apa pun yang berasal dari percakapan seseorang.
\end{icard}
\end{minipage}%
\begin{iband}[signal]
Uji yang harus lulus sebelum tayang: \textbf{buat ingatan sebagai pengguna A, lalu cari sebagai pengguna B.} Kalau kembali, sistemnya belum siap.
\end{iband}
\end{frame}

\framekicker{Risiko}
\begin{frame}[shrink=25]{Data pribadi yang masuk lewat pintu belakang}
{\fontsize{9.4}{12.6}\selectfont\color{fg2}Sistem yang dirancang tanpa data pribadi bisa mengumpulkannya sendiri lewat ingatan: pengguna menyebut nama, nomor, atau kondisi kesehatan di tengah percakapan, dan mekanisme ``ingat hal penting'' menyimpannya dengan patuh.\par}\vskip 4pt
\begin{isteps}
\item \{'h': 'Saring saat menulis, bukan saat membaca', 'p': 'Sekali tersimpan, ia sudah ada di cadangan dan di indeks. Penyaringan di sisi baca datang terlambat.'\}
\item \{'h': 'Simpan kategori, bukan nilainya', 'p': '``pelanggan memilih dihubungi lewat telepon'' cukup; nomor teleponnya tidak perlu ada di ingatan agen.'\}
\item \{'h': 'Perlakukan ingatan sebagai sistem yang didaftarkan', 'p': 'Punya pemilik, punya kebijakan retensi, dan muncul di daftar sistem yang memproses data pribadi --- atau ia tidak boleh ada.'\}
\end{isteps}
\begin{iband}[signal]
Ini kelanjutan langsung dari prinsip di Bab 1: \textbf{batas kemampuan ada di kode.} Ingatan yang tidak bisa menyimpan nomor telepon lebih kuat daripada ingatan yang diminta tidak menyimpannya.
\end{iband}
\end{frame}

\coursesection{06}{Mengukur dan memperbaiki}{Ingatan yang tidak diukur akan membesar sampai jadi masalah.}

\framekicker{Praktik}
\begin{frame}[shrink=25]{Empat angka yang memberi tahu ingatan Anda bermasalah}
\begin{itable}{>{\raggedright\arraybackslash}p{0.2632\textwidth}>{\raggedright\arraybackslash}p{0.3008\textwidth}>{\raggedright\arraybackslash}p{0.376\textwidth}}
\thead{Angka} & \thead{Sehat kalau} & \thead{Gejala kalau tidak} \\ \midrule
Token konteks per giliran & Datar & Naik terus $\rightarrow$ tidak ada yang memangkas \\
Potongan diambil lawan dipakai & Sebagian besar dipakai & Banyak diambil, sedikit dikutip $\rightarrow$ pengambilan terlalu longgar \\
Peringkasan per percakapan & Nol atau satu & Berulang $\rightarrow$ pemangkatan di gambar tadi sedang terjadi \\
Ingatan tersimpan per pengguna & Tumbuh pelan lalu datar & Tumbuh terus $\rightarrow$ tidak ada yang kedaluwarsa \\
\end{itable}
{\fontsize{9.4}{12.6}\selectfont\color{fg2}Tidak satu pun dari empat ini menimbulkan galat. Semuanya bergeser pelan, dan gejalanya di permukaan selalu sama: \textbf{agennya terasa makin bodoh dan makin mahal.}\par}\vskip 4pt
\end{frame}

\framekicker{Praktik}
\begin{frame}[shrink=25]{Mengujinya tanpa model}
{\fontsize{9.4}{12.6}\selectfont\color{fg2}Sebagian besar kegagalan ingatan bisa ditangkap uji biasa, sebab yang diuji adalah \textbf{penyusun konteks}, bukan modelnya.\par}\vskip 4pt
\begin{minipage}[t]{0.494\textwidth}
\begin{icardgood}
\cardh{Uji isi konteks}
Beri riwayat panjang buatan; pastikan batasan dari giliran pertama masih ada di konteks giliran ke-30.
\end{icardgood}
\end{minipage}%
\hfill
\begin{minipage}[t]{0.494\textwidth}
\begin{icardgood}
\cardh{Uji batas pemilik}
Simpan sebagai A, cari sebagai B, harapkan kosong. Ini uji keamanan, dan ia deterministik.
\end{icardgood}
\end{minipage}%


\vskip 4pt

\begin{minipage}[t]{0.494\textwidth}
\begin{icard}
\cardh{Uji pemangkatan}
Ringkas sepuluh kali; periksa apakah angka penting masih ada. Kalau hilang, arsitekturnya yang salah.
\end{icard}
\end{minipage}%
\hfill
\begin{minipage}[t]{0.494\textwidth}
\begin{icard}
\cardh{Uji kutipan}
Pengenal yang dikutip harus ada di hasil pengambilan. Ini pemeriksaan kode, bukan penilaian model.
\end{icard}
\end{minipage}%
\end{frame}

\framekicker{Praktik}
\begin{frame}[shrink=25]{Urutan memperbaikinya, dari yang paling sering berhasil}
\begin{isteps}
\item \{'h': 'Perbaiki apa yang masuk sebelum menambah apa pun', 'p': 'Sebagian besar keluhan \textbackslash\{\}\textbackslash\{\}u201cagennya lupa\textbackslash\{\}\textbackslash\{\}u201d adalah konteks yang salah disusun, bukan ingatan yang kurang.'\}
\item \{'h': 'Simpan aslinya di luar konteks', 'p': 'Ini satu perubahan yang mematahkan pemangkatan peringkasan.'\}
\item \{'h': 'Jadikan pengambilan sebagai alat', 'p': 'Sekaligus memberi Anda jejak dan kutipan --- dua hal yang dibutuhkan bab 7.'\}
\item \{'h': 'Baru tambahkan ingatan jangka panjang', 'p': 'Dengan pemilik, asal-usul, dan kedaluwarsa sejak hari pertama. Menambahkannya belakangan berarti memigrasi data pribadi.'\}
\end{isteps}
\begin{iband}[signal]
Perhatikan bahwa tiga langkah pertama \textbf{tidak menambah komponen baru}. Ingatan jangka panjang adalah yang terakhir dipasang, dan sering ternyata tidak diperlukan.
\end{iband}
\end{frame}

\framekicker{Penutup}
\begin{frame}[shrink=25]{Yang dibawa pulang dari bab ini}
\begin{isteps}
\item \{'h': 'Tidak ada ingatan; yang ada keputusan rekayasa', 'p': 'Dan itu kabar baik: yang Anda susun bisa diperiksa, diuji, dan dipertanggungjawabkan.'\}
\item \{'h': 'Memangkas kehilangan dengan jujur, meringkas tidak', 'p': 'Ringkasan atas ringkasan memangkatkan kerugiannya. Simpan aslinya, ringkas dari sana.'\}
\item \{'h': 'Yang tidak terambil tidak bisa dijawab', 'p': 'Mutu sistem pengambilan ada di pengambilannya, bukan di modelnya.'\}
\item \{'h': 'Konteks adalah spesifikasinya', 'p': 'Dirakit tiap giliran, jadi bisa diuji tiap giliran.'\}
\item \{'h': 'Ingatan adalah basis data baru', 'p': 'Dengan pemilik, retensi, dan cara menghapus --- atau ia jadi kewajiban yang tidak bisa dipenuhi.'\}
\end{isteps}
\end{frame}
\note{Pertanyaan penutup: dari sistem yang kalian bangun, apa yang sebenarnya perlu diingat lintas percakapan? Jawabannya hampir selalu jauh lebih sedikit daripada yang disimpan.}

\end{document}
